
\documentclass[12pt,a4paper]{report}

\usepackage[margin=1in]{geometry}
\usepackage{xcolor}
\usepackage{sectsty}
\usepackage{fancyhdr}
\usepackage{setspace}
\usepackage{amsmath}
\usepackage{amssymb}

\chapterfont{\color{red}\bfseries\Huge}
\sectionfont{\color{magenta}\itshape\Large}
\subsectionfont{\color{orange}\bfseries}

\setstretch{1.3}

\pagestyle{fancy}
\fancyhf{}

\fancyhead[L]{Ketan Sharma}
\fancyhead[C]{Technical Writing Lab}
\fancyhead[R]{\today}

\renewcommand{\headrulewidth}{0.4pt}
\renewcommand{\footrulewidth}{0pt}

\setlength{\headheight}{15pt}

\begin{document}
	
	\begin{titlepage}
		\thispagestyle{empty}
		
		\begin{center}
			
			\vspace*{2cm}
			
			{\Huge\bfseries Introduction to LaTeX}
			
			\vspace{2cm}
			
			{\Large Ketan Sharma}
			
			\vspace{0.5cm}
			
			\textbf{ID:} 2025UCP1398
			
			\vspace{1.5cm}
			
			\textbf{Course:} Technical Writing Lab
			
			\vspace{0.5cm}
			
			\textbf{Course Instructor:} Neeta Nain Ma'am
			
			\vfill
			
			\textbf{Date:} \today
			
		\end{center}
		
	\end{titlepage}
	
	\pagestyle{fancy}
	
	\begin{abstract}
		
		This document provides an introduction to LaTeX, a document preparation system widely used for creating academic papers, technical reports, research articles and presentations. It explains the basic structure of a LaTeX document, formatting of chapters and sections, mathematical expressions, and the use of LaTeX for professional document preparation. It also demonstrates how to create a structured report with customized headings and page headers.
		
	\end{abstract}
	
	\tableofcontents
	
	\chapter{Introduction}
	
	\section{Overview of LaTeX}
	
	LaTeX is a document preparation system based on the TeX typesetting system. It is commonly used in academic and technical fields because it provides precise control over document formatting.
	
	\section{Importance of LaTeX}
	
	LaTeX is useful for preparing documents that contain complex mathematical equations, tables, figures, references and structured content.
	
	\subsection{Applications of LaTeX}
	
	Some common applications include:
	
	\begin{itemize}
		\item Academic research papers
		\item Technical reports
		\item Project documentation
		\item Books and presentations
		\item Mathematical and scientific documents
	\end{itemize}
	
	\chapter{Basic Structure of LaTeX}
	
	\section{Document Structure}
	
	A basic LaTeX document consists of two main parts: the preamble and the document environment.
	
	\begin{verbatim}
		\documentclass{article}
		
		\begin{document}
			
			Hello World!
			
		\end{document}
	\end{verbatim}
	
	The preamble contains document settings and packages, while the document environment contains the actual content.
	
	\section{Formatting Text}
	
	LaTeX provides several commands for formatting text.
	
	\subsection{Text Formatting Examples}
	
	\begin{itemize}
		\item \textbf{Bold Text}
		\item \textit{Italic Text}
		\item \underline{Underlined Text}
	\end{itemize}
	
	\chapter{Mathematics in LaTeX}
	
	\section{Mathematical Expressions}
	
	LaTeX provides powerful mathematical typesetting capabilities.
	
	An example of an inline equation is $E=mc^2$.
	
	A displayed mathematical equation is:
	
	\[
	a^2+b^2=c^2
	\]
	
	\section{Fractions and Summations}
	
	A fraction can be written as:
	
	\[
	\frac{a+b}{c+d}
	\]
	
	A summation can be written as:
	
	\[
	\sum_{i=1}^{n} i = \frac{n(n+1)}{2}
	\]
	
	\subsection{Matrices}
	
	Matrices can also be created using LaTeX:
	
	\[
	A=
	\begin{bmatrix}
		1 & 2\\
		3 & 4
	\end{bmatrix}
	\]
	
	\chapter{Conclusion}
	
	LaTeX is a powerful tool for preparing professional documents, academic reports and presentations. Its structured approach makes it suitable for handling complex mathematical expressions, tables, references and formatting requirements.
	
	Learning LaTeX helps students create well-organized and consistent technical documents.
	
\end{document}
